\documentclass{article}
\usepackage{amsmath, amssymb, amsthm}
\usepackage{hyperref}
\usepackage{geometry}
\geometry{margin=1in}

\title{Erd\H{o}s Problem \#301}
\date{Last updated: 2025-08-31}
\author{Source: erdosproblems.com}

\begin{document}
\maketitle

\noindent\textbf{Status:} OPEN \quad \textbf{No prize}

\noindent\textbf{Tags:} number theory, unit fractions

\medskip
\noindent\textbf{Problem Statement:}

Let $f(N)$ be the size of the largest $A\subseteq \{1,\ldots,N\}$ such that there are no solutions to\[\frac{1}{a}= \frac{1}{b_1}+\cdots+\frac{1}{b_k}\]with distinct $a,b_1,\ldots,b_k\in A$?

Estimate $f(N)$. In particular, is it true that $f(N)=(\tfrac{1}{2}+o(1))N$?

\bigskip
\noindent\textbf{Remarks:}

The example $A=(N/2,N]\cap \mathbb{N}$ shows that $f(N)\geq N/2$.

Wouter van Doorn has given an elementary argument that proves\[f(N)\leq (25/28+o(1))N.\]Indeed, consider the sets $S_a=\{2a,3a,4a,6a,12a\}\cap [1,N]$ as $a$ ranges over all integers of the form $8^b9^cd$ with $(d,6)=1$. All such $S_a$ are disjoint and, if $A$ has no solutions to the given equation, then $A$ must omit at least two elements of $S_a$ when $a\leq N/12$ and at least one element of $S_a$ when $N/12<a\leq N/6$, and an elementary calculation concludes the proof.

Stijn Cambie and Wouter van Doorn have noted that, if we allow solutions to this equation with non-distinct $b_i$, then the size of the maximal set is at most $N/2$. Indeed, this is the classical threshold for the existence of some distinct $a,b\in A$ such that $a\mid b$.

See also [302] and [327].

\bigskip
\noindent\small{Source: \url{https://www.erdosproblems.com/301}}
\end{document}
